%%%%%%%%%%%%%%%%%%%%%%%%%%%%%%%%%%%%%%%%%%%%%%%%%%%%%%%%%%%%%%%%%%%%%%
% Lớp 10 | Đại số tổ hợp | Quy tắc cộng-quy tắc nhân
% Dạng 0D8?1-6: Sơ đồ hình cây
%%%%%%%%%%%%%%%%%%%%%%%%%%%%%%%%%%%%%%%%%%%%%%%%%%%%%%%%%%%%%%%%%%%%%%


% nguồn: TuLuan_Lop10_Dot2 (nhập 14/11/2026)
\begin{bt}%[0D8H1-6]%[Dự án đề kiểm tra HKII NH22-23-BCTuan]%[THPT Thực Hành Sài Gòn]%[ma:0fee19fb]
	\begin{enumerate}
		\item Có ba cái hộp, hộp thứ nhất chứa $2$ quả cầu dán nhãn $A$, $B$; hộp thứ hai chứa $3$ quả cầu dán nhãn $a$, $b$, $c$; hộp thứ ba có $2$ quả cầu dán nhãn $1$, $2$. Lấy ngẫu nhiên một quả cầu trong mỗi hộp trên. Hãy vẽ sơ đồ hình cây để thể hiện tất cả các kết quả có thể xảy ra.
		\item Từ các chữ số $1; 2; 3; 4; 5; 6$ có thể lập được bao nhiêu số tự nhiên có bốn chữ số khác nhau và lớn hơn $4500$?
	\end{enumerate}
\loigiai{
\begin{enumerate}
	\item Sơ đồ hình cây
		\begin{center}
		\begin{tikzpicture}[->,>=stealth]
			\tikzstyle{hollow node}=[circle,draw,inner sep=1.5]
			\tikzstyle{solid node}=[circle,draw,inner sep=1.5,fill=black]
			\tikzset{
				red node/.style={circle,draw=blue,inner sep=2.5},
				blue node/.style={circle,draw=blue,inner sep=2.5}
			}
			\tikzstyle{level 1}=[sibling distance=80mm]
			\tikzstyle{level 2}=[sibling distance=25mm]
			\tikzstyle{level 3}=[sibling distance=10mm]
			\node[hollow node]{Gốc}
			child{node[blue node]{A}
				child{node[red node]{a}child{node[blue node]{1}}child{node[blue node]{2}}}
				child{node[red node]{b}child{node[blue node]{1}}child{node[blue node]{2}}}
				child{node[red node]{c}child{node[blue node]{1}}child{node[blue node]{2}}}
			}
			child{node[blue node]{B}
				child{node[red node]{a}child{node[blue node]{1}}child{node[blue node]{2}}}
				child{node[red node]{b}child{node[blue node]{1}}child{node[blue node]{2}}}
				child{node[red node]{c}child{node[blue node]{1}}child{node[blue node]{2}}}
			}
			;
			\path (-8.5,-1.5) node{Hộp I}
			(-8.5,-2.8) node{Hộp II}
			(-8.5,-4.5) node{Hộp III};
		\end{tikzpicture}
	\end{center}
\item Giả sử số cần tìm có dạng $\overline{abcd}$ ($a\ne 0$).\\
Vì số lớn hơn $4500$ nên có các trường hợp sau:\\
{\bf TH1:} $a>4$, khi đó số cách chọn chữ số $a$ là $2$ cách.\\
Số cách chọn $\overline{bcd}$: $\mathrm{A}_5^3$ cách.\\
Vậy trường hợp này có $2\cdot \mathrm{A}_5^3=120$ (cách).\\
{\bf TH2:} $a=4$, khi đó chữ số $b$ có 2 cách chọn ($5$ hoặc $6$).\\
Số cách chọn $\overline{cd}$ là $\mathrm{A}_4^2$ cách.\\
Vậy trường hợp này có $2\cdot \mathrm{A}_4^2=24$ (cách).\\
Theo quy tắc cộng, có $120+24=144$ (số).
\end{enumerate}
}
\end{bt}

% nguồn: TuLuan_Lop10_Dot2 (nhập 14/11/2026)
\begin{bt}%[0D8H1-6]%[Tex hóa SGK CD-CT,T8/22, Le Vu Hai]%[ma:d583f476]
	Bốn bạn Ánh, Bình, Cường, Hoa cùng thi vào trường trung học phổ thông chất lượng cao Bình Minh. Kết quả thi cho bởi bảng thống kê sau$\colon$
	\begin{center}
		\begin{tabular}{|c|c|c|c|}
			\hline Học sinh & Điểm Toán & Điểm Ngữ văn & Điểm Tiếng Anh \\
			\hline Ánh & $10$ & $8$ & $10$ \\
			\hline Bình & $6$ & $7$ & $5$ \\
			\hline Cường & $9$ & $8$ & $9$ \\
			\hline Hoa & $8$ & $6$ & $9$ \\
			\hline
		\end{tabular}
	\end{center}
	Quy định$\colon$
	\begin{itemize}
		\item Điểm trung bình viết tắt là ĐTB và tính theo công thức$\colon$
		$$ \text{ĐTB} = (\text{Điểm Toán} + \text{Điểm Ngữ văn} + \text{Điểm Tiếng Anh}) \colon 3.$$
		\item Điểm trung bình Toán viết tắt là ĐTB $\mathrm{T}$ và tính theo công thức$\colon$
		$$ \text{ĐTB T}=(2 \times \text{Điểm Toán} + \text{Điểm Ngữ văn} + \text{Điểm Tiếng Anh})\colon 4.$$
		\item Điểm trung bình Ngữ văn viết tắt là ĐTB V và tính theo công thức$\colon$
		$$ \text{ĐTB V} = ( \text{Điểm Toán} +2 \times  \text{Điểm Ngữ văn} + \text{Điểm Tiếng Anh}) \colon 4.$$
		\item Điểm trung bình Tiếng Anh viết tắt là ĐTB A và tính theo công thức$\colon$
		$$ \text{ĐTB A}= (\text{Điểm Toán} + \text{Điểm Ngữ văn} +2 \times  \text{Điểm Tiếng Anh}) \colon4.$$
	\end{itemize}
	Cách xét tuyển như sau:
	\begin{itemize}
		\item Trúng tuyển: $\text{ĐTB}>7,5$;
		\item Trúng tuyển lớp Toán: $\text{ĐTB T}>9$;
		\item Trúng tuyển lớp Văn: $\text{ĐTB V} > 8$ ;
		\item Trúng tuyển lớp Tiếng Anh: $\text{ĐTB A}>8,5$.
	\end{itemize}
	Hãy lập bảng kết quả theo mẫu dưới đây và xem mỗi bạn có những cơ hội nào (làm tròn kết quả đến hàng phần trăm).
	\begin{center}
		\begin{tabular}{|c|c|c|c|c|c|}
			\hline Học sinh & ĐTB & ĐTB T & ĐTB V & ĐTB A & Kết quả \\
			\hline Ánh & $?$ & $?$ & $?$ & $?$ & $?$ \\
			\hline Bình & $?$ & $?$ & $?$ & $?$ & $?$ \\
			\hline Cường & $?$ & $?$ & $?$ & $?$ & $?$ \\
			\hline Hoa & $?$ & $?$ & $?$ & $?$ & $?$ \\
			\hline
		\end{tabular}
	\end{center}
	\loigiai{
		\begin{center}
			\begin{tabular}{|c|c|c|c|c|c|}
				\hline Học sinh & ĐTB & ĐTB T & ĐTB V & ĐTB A & Kết quả \\
				\hline Ánh & $9{,}33$ & $9{,}5$ & $9$ & $9{,}5$ & Đậu tất cả các lớp\\
				\hline Bình & $6$ & $6$ & $6{,}25$ & $5{,}75$ & Không đậu \\
				\hline Cường & $8{,}67$ & $8{,}75$ & $8{,}5$ & $8{,}75$ & Đậu lớp V, A \\
				\hline Hoa & $7{,}67$ & $7{,}75$ & $7{,}25$ & $8$ & Đủ điểm trung bình nhưng không đậu lớp T, V, A \\
				\hline
			\end{tabular}
		\end{center}	
	}
\end{bt}
